% ---
% Pacotes fundamentais 
% ---
\usepackage{cmap}				% Mapear caracteres especiais no PDF
\usepackage{times}			    % Usa a fonte Times
\usepackage[T1]{fontenc}		% Selecao de codigos de fonte.fontenc
\usepackage[utf8]{inputenc}		% Codificacao do documento (conversão automática dos acentos)
\usepackage{lastpage}			% Usado pela Ficha catalográfica
\usepackage{indentfirst}		% Indenta o primeiro parágrafo de cada seção.
\usepackage{color}				% Controle das cores
\usepackage{graphicx}			% Inclusão de gráficos
\usepackage{amsfonts}			% Símbolos
%----Ajuste no alinhamento das listas
\usepackage{enumitem}
\setitemize[0]{itemindent=0.4cm,itemsep=0pt}
\setenumerate[0]{itemindent=0.5cm,itemsep=0pt}
%------
% ---
% Pacotes de citações
% ---
\usepackage[brazilian,hyperpageref]{backref}	 					  % Paginas com as citações na bibl

\usepackage[a4paper]{geometry}
%Referência
\usepackage[alf, 	
			 		abnt-emphasize=bf,
				    abnt-url-package=none,
				    abnt-repeated-title-omit=yes,
				    abnt-full-initials=yes,                                        %yes nome por extenso, no apenas iniciais
					abnt-etal-list=3												%abreviar com mais de 3 autores
]{abntex2/abntex2cite}				 														    % Citações padrão ABNT
\usepackage{lipsum}							   								       % para geração de dummy text
\usepackage{multirow}
\usepackage{float}
\usepackage{setspace}
\usepackage{tabularx}
\usepackage[table,xcdraw]{xcolor}
\usepackage{booktabs}			% \toprule, \midrule e \bottomrule usados nas tabelas do modelo
\usepackage{longtable}
\usepackage{array}
\newcolumntype{L}[1]{>{\raggedright\arraybackslash}p{#1}}

% --- Figuras e gráficos vetoriais
\usepackage{tikz}
\usetikzlibrary{positioning,arrows.meta,shapes.geometric,shapes.multipart,fit,backgrounds,calc,shadows}
\usepackage{pgfplots}
\pgfplotsset{compat=1.18, /pgf/number format/use comma, /pgf/number format/1000 sep={.}}
\renewcommand{\arraystretch}{1.2}

\tikzset{
  caixa/.style={draw, rounded corners=2pt, align=center, font=\footnotesize, minimum height=0.9cm, inner sep=4pt, fill=white},
  ia/.style={caixa, fill=black!8},
  ext/.style={caixa, dashed},
  seta/.style={-{Stealth[length=2.2mm]}, thick},
  grupo/.style={draw, rounded corners=4pt, inner sep=8pt, dashed, black!60},
  rotulo/.style={font=\scriptsize\itshape, align=center}
}

% --- Código-fonte nos apêndices
\usepackage{listings}
\lstdefinelanguage{JavaScript}{
  keywords={const,let,var,function,return,if,else,for,of,new,true,false,null,typeof,throw,try,catch},
  sensitive=true,
  comment=[l]{//},
  morecomment=[s]{/*}{*/},
  morestring=[b]',
  morestring=[b]",
  morestring=[b]`
}
\lstset{
  basicstyle=\ttfamily\scriptsize,
  breaklines=true,
  columns=fullflexible,
  keepspaces=true,
  frame=single,
  numbers=left,
  numberstyle=\tiny,
  showstringspaces=false,
  keywordstyle=\bfseries,
  commentstyle=\itshape\color{black!60},
  literate={á}{{\'a}}1 {à}{{\`a}}1 {ã}{{\~a}}1 {â}{{\^a}}1 {é}{{\'e}}1 {ê}{{\^e}}1 {í}{{\'i}}1
           {ó}{{\'o}}1 {õ}{{\~o}}1 {ô}{{\^o}}1 {ú}{{\'u}}1 {ç}{{\c c}}1 {Á}{{\'A}}1 {É}{{\'E}}1
           {Í}{{\'I}}1 {Ó}{{\'O}}1 {Ú}{{\'U}}1 {Ç}{{\c C}}1 {Ã}{{\~A}}1 {—}{{---}}1 {–}{{--}}1
           {⚙}{{}}0 {·}{{$\cdot$}}1 {Õ}{{\~O}}1 {Ê}{{\^E}}1 {️}{{}}0 {→}{{$\rightarrow$}}1
}
%\usepackage[none]{hyphenat}
%\captionsetup[table]{justification=raggedright}
% Configurações de aparência do PDF final
% alterando o aspecto da cor azul
\definecolor{blue}{RGB}{41,5,195}

% --- 
% Espaçamentos entre linhas e parágrafos 
% --- 
% O tamanho do parágrafo é dado por:
\setlength{\parindent}{2cm}
\linespread{1.5}

%Espaçamento depois dos títulos
\setlength{\afterchapskip}{\baselineskip}
% %\setlength{\afterchapskip}{\lineskip}

% Controle do espaçamento entre um parágrafo e outro:
\setlength{\parskip}{0cm}  % tente também \onelineskip
\setlength{\emergencystretch}{2em} % evita linhas além da margem em parágrafos com termos longos

\hangcaption
\captionstyle[\raggedright]{}

%Estava mostrando nas referencias quais paginas estavam sendo referenciadas
\renewcommand{\backref}{}
\renewcommand*{\backrefalt}[4]{}

%Reduzir a fonte do caption
%\captionnamefont{\centering\ABNTEXfontereduzida}
%\captiontitlefont{\centering\ABNTEXfontereduzida}
%Ajuste nas listas de tabela, ilustrações e quadros
\setlength\cftbeforechapterskip{0pt}
% ---
% compila o indice
% ---
\makeindex
% ---